% ============================================================
% pretty-charts TikZ 时序图预置样式
% 依赖：先 \input preamble.tex 与 flowchart-styles.tex（复用其 \pcPaperMode/\pcPPTMode 模式变量）
% 论文版：白底黑框细线；演示版：彩色粗线（随模式自动切换）
% ============================================================
\usetikzlibrary{calc}

\tikzset{
  % 参与者（顶部矩形；角色栏使用同一样式）
  seq participant/.style={draw=\pcNodeLine, fill=\pcNodeFill, minimum width=1.9cm,
                          minimum height=0.65cm, font=\pcNodeFont, text=\pcNodeText, align=center},
  % 生命线：细虚线
  seq lifeline/.style={draw=black!55, line width=0.4pt, dash pattern=on 2pt off 2pt},
  % 消息（同步实线 / 返回虚线）
  seq msg/.style={-{Stealth[length=2.2mm]}, line width=0.5pt, draw=\pcArrow},
  seq ret/.style={seq msg, dashed},
  % 消息文字与序号
  seq label/.style={font=\footnotesize, text=\pcNodeText, inner sep=2pt, fill=white},
  seq num/.style={circle, fill=black!80, text=white, font=\tiny\bfseries, inner sep=1.2pt},
  % 激活条（可选）
  seq activation/.style={draw=black!60, fill=black!4, line width=0.4pt, rounded corners=1pt},
  % alt/loop 片段框：复用流程图阶段容器
  seq frame/.style={pc sub},
  seq frame label/.style={pc sub label, anchor=north west},
}

% ---- 用法示例 ----
% \node[seq participant] (u) at (0,0) {用户};
% \draw[seq lifeline] (u.south) -- ++(0,-8);
% \draw[seq msg] ($(u.south)+(0,-1)$) -- node[seq label, above=1pt]{请求} ($(g.south)+(0,-1)$);
% \node[seq num] at ($(u.south)+(0,-1)+(0.2,0.22)$) {1};
